\subsection{Limpeza e transformação}
A fonte foi preservada. Padronizamos grafias, separadores decimais e códigos de ausência, convertendo as medidas às unidades do dicionário. As regras por registro usam coerência entre campos e limites amplos de plausibilidade. Medidas suspeitas são mascaradas com registro da regra. Conservar a linha permite usar suas demais informações. A correção de altura em centímetros depende da coerência com a escala do campo. Para IMC, usamos a relação peso dividido pela altura em metros ao quadrado quando ambos os insumos são válidos, com preservação controlada do valor histórico quando essa reconstrução é inviável. Para IDADE, a diferença entre nascimento e atendimento é usada quando a interpretação das datas é coerente. O valor informado tem papel de referência nas ambiguidades de formato e século. As regras e exceções por variável foram documentadas na auditoria.

Em cada treino, comparamos a mediana global com medianas em faixas fixas de idade para Peso, Altura, IMC e FC quando presentes. A faixa requer pelo menos 30 valores observados no treino. Caso contrário ou com idade indisponível, vale a mediana global. O indicador de imputação numérica conserva informação sobre indisponibilidade e o ensaio de máscaras investiga seu poder preditivo. SEXO usa codificação \textit{one-hot} com vocabulário do treino. Para a logística comparamos padronização das medidas originais com a padronização após aplicar \(\log(1+\text{Peso})\) ao Peso. As árvores usam as unidades originais. Indicadores binários permanecem em sua escala. Nenhum parâmetro de transformação é estimado nas linhas de validação.

\subsection{Mineração e avaliação}
O cenário \fisio{} usa Peso, Altura, IMC, IDADE, pressões sistólica e diastólica, FC e SEXO. \strict{} usa Peso, Altura, IDADE e SEXO. A segunda análise reduz a dependência de medidas cardiovasculares cuja relação com a rotulagem é incerta e elimina a redundância de IMC com os insumos antropométricos. Os dois cenários compartilham os mesmos registros e cortes externos.

Usamos a classe majoritária como referência, uma árvore rasa como comparação interpretável, regressão logística regularizada e floresta aleatória \cite{breiman2001,pedregosa2011}. A grade por cenário inclui 105 configurações. São 48 de árvore (critério, profundidade, tamanho mínimo de folha, pesos e imputação), 32 de logística (regularização, representação, pesos e imputação), 24 de floresta (profundidade, folha mínima, pesos e imputação, com 150 árvores) e uma referência majoritária. Os pesos balanceados aumentam a contribuição da classe menos frequente no treino. A busca compara seu efeito com pesos uniformes. O limiar de 0,5 é fixado para a comparação.

A validação cruzada aninhada separa ajuste e estimativa de desempenho \cite{varma2006}. Em cada um dos cinco folds externos, três folds internos escolhem uma configuração por família pela média da acurácia balanceada (BA). Em empate numérico, prevalece a ordem fixa da grade. O procedimento reajusta o vencedor no treino externo e gera uma previsão para cada linha de sua validação. A estratificação procura preservar as proporções de classe e o agrupamento evita separar componentes candidatos entre treino e validação. Usamos \texttt{random\_state=42} nos estimadores estocásticos, semente 743 nos cortes externos e \(811+k\) no corte interno do fold \(k\), com \(k=0,\ldots,4\).

A BA é a média da sensibilidade (fração de \textit{Anormal} recuperada) e da especificidade (fração de \textit{Normal} reconhecida). Ela dá o mesmo peso às classes. Relatamos também precisão da classe positiva, \(F_1\), precisão média (AP) e matriz de confusão. As métricas principais reúnem as previsões externas de todos os folds. O desvio-padrão entre folds descreve a variação observada. Como os treinos se sobrepõem, esse desvio não foi usado como erro-padrão de observações independentes.
